\section{A direct equation in 83 unrestricted integers}\label{sec:direct-results}

The modular witness-reuse theorem leads to a concrete second construction.
Its source is the unmodified direct numerical-calculus polynomial in
package 166, followed by the two semantic changes just proved and explicit
integral coordinate elimination. The main endpoint is
\[
 D_{83}(x_0,\ldots,x_{82})
       =\sum_{j=0}^{30}R_j(x_0,\ldots,x_{82})^2.
 \tag{Direct}
\]
Appendix~\ref{app:direct83-definition} prints every free assignment, every
ordered substitution, and every residual. The electronic appendix also
supplies the corresponding literal arithmetic circuit and source coefficients.
All \(83\) inputs range independently over \(\Z\).
There are no hidden inequalities or extra witnesses.

\begin{theorem}[Direct C222 route]\label{thm:direct}
For the fixed polynomial (Direct), with the inherited C222/ZF/ZFC
refutation calibration identified in Section~\ref{sec:boundary},
\[
 \ZFC\vdash
 \left(\neg\exists\boldsymbol x\in\Z^{83}\;D_{83}(\boldsymbol x)=0\right)
       \longleftrightarrow\Con(\ZFC).
\]
The exact total degree is \(624\), and the preferred factored evaluator has
\(1{,}789\) live binary arithmetic operations. Its integer literal scalars
have absolute value at most \(64\).
\end{theorem}

This route does not simulate K215 or specialize Jones's universal
polynomial. It retains a specific direct proof calculus and its previously
identified consistency calibration. Compared with the \(2{,}428\)-operation \(92\)-integer, degree-\(692\) equation
\cite{source175}, the new point improves dimension \(92\to83\), degree
\(692\to624\), and operation count \(2{,}428\to1{,}789\) together.
The latter count uses the exact factorizations described below; the
canonical evaluator retained for the projection audit has \(2{,}939\) operations.

\subsection{A signed frame lowers the source degree}

For the source used in this endpoint, one may replace \(q^2W_0\) by
\(qW_0\). During completeness, the semantic construction gives
\(q=(BH)^x\), with \(x\ge2\), and every prime dividing the odd modulus
\(\mathcal M=Nq^x\) divides \(q\). Thus every such prime has square
dividing \(W=qW_0\). The proof of Theorem~\ref{thm:reuse} only needs
that local square divisibility, so the same modular construction works.
The positivity of the semantic \(q\) gives a positive completeness frame.

Soundness must also cover a negative value of the unclassified input \(q\).
The affine source row gives \(|q|\ge3\), so \(F=|W|\ge3W_0\).
The new \(A_*=L_0^2+2r^2\) is positive and at least \(L_0^2\).
When \(W=-F\), the frame equation becomes
\[
 L_0^2-F^3(F-2)r^2=1.
\]
This is the Pell equation at parameter \(F-1\), with ordinate \(F|r|\).
The congruence
\[
 \PellY_j(F-1)\equiv(-1)^{j-1}j\pmod F
\]
forces its positive classified index to be a multiple of \(F\).
As \(r\) is odd, it is nonzero, so that index is at least \(F\).
Consequently
\[
 A_*\ge L_0^2\ge(F-1)^{2F}>F^F\qquad(F\ge3).
\]
For \(W>0\) the retained positive-frame argument gives the required
bound. Every old bounded atom has absolute value below \(F\).
The source's only uses of \(W\) and its private frame wires are the frame
equations and its definition, checked by an occurrence audit.
The inherited signed power recognition therefore applies with \(F=|W|\).

The signed-frame row has degree three, rather than four.
Replaying the inherited integral coordinate changes, with the two private
center inputs deleted, gives an explicit equation in \(227\) integers,
the sum of \(175\) residual squares, of exact degree six.
Its collected polynomial has \(32{,}328\) monomials and height \(221{,}184\).
The archive also retains the positive-frame degree-eight alternative.
The \(227\)-variable intermediate is a proof source, not a global quartic
or low-dimensional record.

\subsection{Acyclic unit elimination}

Let \(S_0,\ldots,S_{174}\) be the literal residuals in that \(227\)-variable
source. A selected row has the form
\[
 S_r=\epsilon x_v+H,\qquad \epsilon\in\{1,-1\},
 \qquad H\in\Z[x_0,\ldots,\widehat{x_v},\ldots,x_{226}].
 \tag{Pivot}
\]
Set \(x_v=-\epsilon H\) and remove that row.
The dependence graph of the selected definitions is acyclic, so the
removed coordinates are restored as integer polynomials in the retained
ones, in a finite order. No division by a nonunit is performed.

For the principal endpoint the certificate contains \(144\) such
definitions, leaving \(83\) coordinates and \(31\) residuals.
If \(\Phi:\Z^{83}\to\Z^{227}\) is the resulting reconstruction, exact
polynomial substitution gives
\[
 S_r(\Phi(\boldsymbol x))=0
 \quad\text{for every removed row},\qquad
 D_{227}(\Phi(\boldsymbol x))=D_{83}(\boldsymbol x).
 \tag{Restore}
\]
Projection and reconstruction are inverse on the integer zero sets.
The signed-frame and witness-reuse changes earlier preserve existence;
the acyclic elimination itself gives this stronger zero-set bijection.
The method of unit elimination is inherited from the repository's prior
projection work. The new semantic source, selected dependency orientation,
literal certificates, and measured endpoints are supplied here.

\subsection{Exact degree and cost certificates}

The circuit gives an upper bound \(312\) on each retained residual degree.
A univariate specialization \(x_i=d_i t\), with explicitly supplied
integer \(d_i\), has a nonzero \(t^{312}\) coefficient in a retained
residual modulo \(1{,}000{,}003\). This proves that residual has exact
degree \(312\) over \(\Z\).
The sum of squares has degree \(624\): highest real homogeneous squares
cannot cancel. Independent finite-field nonzero partial derivatives
certify that every one of the \(83\) declared inputs actually occurs.

Exact propagation of coefficient-norm upper bounds through the circuit
gives the rigorous bound
\[
 \max |\text{expanded coefficient of }D_{83}|<2^{1568}.
\]
This is a bound, not an exact expanded height. A literal scalar bound
of \(64\) in the factored evaluator is a different measure. The complete expanded degree-\(624\)
polynomial was not collected; the circuit is its exact finite definition.

\subsection{A factored evaluator for the same polynomial}
\label{sec:factored-evaluator}

The expanded source-coordinate presentation contains two long expressions
that share a quadratic form. Factoring these expressions changes neither
the source, the projection, nor the polynomial. It gives a second literal
evaluator with \textbf{1,789 live binary operations}, comprising 955
multiplications and 834 additions. The largest absolute literal is 64.
The canonical evaluator has 2,939 operations and largest literal 324.
Both define exactly the same polynomial in 83 unrestricted integers, of
degree 624, with 31 residual squares and expanded coefficient height below
$2^{1568}$.

Use the source coordinates $z_0,\ldots,z_{226}$ from the projection appendix,
and introduce the following polynomial abbreviations:
\begin{align*}
 H_*={}&z_8+z_9+4(z_{19}+z_{24}+z_{30}+z_{36}+z_{65})\\
       &+2(z_{22}+z_{27}+z_{29}+z_{32}+z_{34}
                    +z_{37}+z_{40}+z_{44}+z_{66}),\\
 q_*={}&3+6z_{220},\\
 \mathcal I={}&\{1,2,10,11,12,13,14,15,16,17,103,121,127,195\},\\
 \mathcal J={}&\{145,153,158,161,164,167,170,173,
                  177,180,183,187,190,214,215,216\},\\
 Q_*={}&z_0^2+(z_0+z_4)^2+\sum_{j\in\mathcal I}z_j^2,\\
 N_*={}&z_{193}^2+2z_{158}-z_{216}-\sum_{j\in\mathcal J}z_j^2.
\end{align*}
The long restoration of $z_{139}$ has the exact factorization
\begin{align}
 z_{139}={}&
 q_*\bigl[Q_*-2z_8+4z_{225}(z_{225}+1)+q_*^2+4\bigr]\notag\\
 &+2(1+3z_{220})(H_*^2+z_8^2+z_9^2)
       -2z_8z_9+z_{19}+N_*-3.\label{eq:factored-z139}
\end{align}
The corresponding retained residual is
\begin{equation}
 R_{28}=N_*-1-z_{19}-H_*^2+(z_8+z_9)^2.
 \label{eq:factored-R28}
\end{equation}
Here residual indexing is zero based: $R_{28}$ is retained source row 170.
These letters denote expression abbreviations, not additional unknowns.

The independent checker expands~\eqref{eq:factored-z139} and
\eqref{eq:factored-R28} in the free polynomial ring
$\mathbb Z[z_0,\ldots,z_{226}]$ and compares all 336 and 154 nonzero
coefficients, respectively, with the frozen definitions. It then checks
every restoration, every retained residual, and the sum-of-squares output
in both complete arithmetic DAGs. Thus the equality holds at every integer
tuple, not only on their zero sets. The original degree and occurrence
certificates also replay for the alternate evaluator. Exact propagation
of coefficient $\ell_1$ majorants again gives the strict bound $2^{1568}$;
this is distinct from the largest literal 64.

The alternate graph, formulas, byte-reproducible builder, independent
standard-library checker, and receipt are in
\texttt{direct/factored83/}. The canonical graph and certificates remain
in \texttt{direct/P83/}. No semantic theorem is changed by this evaluator
optimization.


\begin{center}
\begin{tabular}{rrrr}
\toprule
Integer variables & Exact degree & Residual squares & Interpretation\\
\midrule
79 & 8,388,620 & 27 & Fewest coordinates in this family\\
81 & 2,048 & 29 & Smaller dimension, higher degree\\
82 & 764 & 30 & Intermediate point\\
\textbf{83} & \textbf{624} & \textbf{31} & \textbf{Principal direct endpoint}\\
87 & 268 & 35 & Lower-degree alternative\\
89 & 128 & 37 & Lower-degree alternative\\
91 & 80 & 39 & Lower-degree alternative\\
\bottomrule
\end{tabular}
\end{center}

Each row is an independently specified ordinary polynomial over unrestricted
integers, with its own literal circuit and exact degree certificate.
The \(79\)-variable equation improves on the \(87\)-variable coordinate
count in the inspected oriented-power family, at a slightly higher
degree than its \(8{,}388{,}612\) endpoint \cite{source177}.
The table makes the degree cost explicit. It is not a claim that an
equation in \(79\) variables is simpler than every equation in fewer
variables using the large native coefficient tuple.
